% GENERATED FILE. Edit canonical lessons, then run npm run generate:books.
% canonical-source-hash: fnv1a64:8a39a70db0d7d5c8
% canonical-lessons: JA-C477-oya, JA-C477-oyako, JA-C477-suekko, JA-C477-mendou, JA-C477-shitsuke

\chapter{Parents and Children}
\label{ch:ja-parents-and-children}
\hlchaptermodality{477}

\begin{chapteropening}
\textbf{By the end of this chapter:} \emph{I can talk about parents and children.}

Complete the last lesson of chapter 477: I can talk about parents and children.
\end{chapteropening}

\section[\texorpdfstring{\ja{おや} (oya)}{おや (oya)}]{\ja{おや} (oya) — a parent}

\label{lesson:JA-C477-oya}



\begin{quote}
Before the new one: say the Japanese for the eldest son, then the Japanese for the eldest daughter.
\end{quote}

\subsection*{You'll want to know: \ja{おや}}
\textbf{\ja{おや}} — \emph{oya} — \textquotedblleft{}a parent\textquotedblright{}.

\begin{quote}
\textbf{\ja{おやは} \ja{こどもを} \ja{しんぱいします。}} — \emph{oya wa kodomo o shinpai shimasu} — \textquotedblleft{}Parents worry about their children.\textquotedblright{}
\end{quote}

\begin{cousinweb}
\textbf{\ja{おや}} (\emph{oya}) is a native Japanese word (wago), inherited rather than borrowed from Chinese; its kanji spelling was attached to it later (JMdict).
\end{cousinweb}

\begin{grammarlens}[title={What you've built}]
One more word for describing what you lived and felt.
\end{grammarlens}

\subsection*{Your turn}
\begin{itemize}
\raggedright
  \item \emph{Say it:} \emph{oya}
  \item \emph{Say it:} \emph{oya}, once more
  \item \emph{Say it:} \emph{oya wa kodomo o shinpai shimasu} — Parents worry about their children.
  \item \emph{Recall:} say the Japanese for the eldest son, then the Japanese for the eldest daughter, then say \emph{oya} again
\end{itemize}

\begin{tcolorbox}[breakable,colback=teal!4,colframe=teal!35!black,title={Before you move on}]
What does \ja{おや} mean? (\textquotedblleft{}A parent\textquotedblright{}.) Say it once more.
\end{tcolorbox}
\section[\texorpdfstring{\ja{おやこ} (oyako)}{おやこ (oyako)}]{\ja{おやこ} (oyako) — parent and child}

\label{lesson:JA-C477-oyako}



\begin{quote}
Before the new one: say the Japanese for the eldest daughter, then the Japanese for a parent.
\end{quote}

\subsection*{You'll want to know: \ja{おやこ}}
\textbf{\ja{おやこ}} — \emph{oyako} — \textquotedblleft{}parent and child\textquotedblright{}.

\begin{quote}
\textbf{\ja{こうえんで} \ja{おやこが} \ja{あそんでいます。}} — \emph{kōen de oyako ga asonde imasu} — \textquotedblleft{}A parent and child are playing in the park.\textquotedblright{}
\end{quote}

\begin{cousinweb}
\textbf{\ja{おやこ}} (\emph{oyako}) is a native Japanese compound (wago), inherited rather than borrowed from Chinese: \emph{oya}, \textquotedblleft{}parent\textquotedblright{} and \emph{ko}, \textquotedblleft{}child\textquotedblright{} (JMdict, KANJIDIC2).
\end{cousinweb}

\begin{grammarlens}[title={What you've built}]
One more word for describing what you lived and felt.
\end{grammarlens}

\subsection*{Your turn}
\begin{itemize}
\raggedright
  \item \emph{Say it:} \emph{oyako}
  \item \emph{Say it:} \emph{oyako}, once more
  \item \emph{Say it:} \emph{kōen de oyako ga asonde imasu} — A parent and child are playing in the park.
  \item \emph{Recall:} say the Japanese for the eldest daughter, then the Japanese for a parent, then say \emph{oyako} again
\end{itemize}

\begin{tcolorbox}[breakable,colback=teal!4,colframe=teal!35!black,title={Before you move on}]
What does \ja{おやこ} mean? (\textquotedblleft{}Parent and child\textquotedblright{}.) Say it once more.
\end{tcolorbox}
\section[\texorpdfstring{\ja{すえっこ} (suekko)}{すえっこ (suekko)}]{\ja{すえっこ} (suekko) — the youngest child}

\label{lesson:JA-C477-suekko}



\begin{quote}
Before the new one: say the Japanese for a parent, then the Japanese for parent and child.
\end{quote}

\subsection*{You'll want to know: \ja{すえっこ}}
\textbf{\ja{すえっこ}} — \emph{suekko} — \textquotedblleft{}the youngest child\textquotedblright{}.

\begin{quote}
\textbf{\ja{すえっこは} \ja{まだ} \ja{ちいさいです。}} — \emph{suekko wa mada chiisai desu} — \textquotedblleft{}The youngest child is still small.\textquotedblright{}
\end{quote}

\begin{cousinweb}
\textbf{\ja{すえっこ}} (\emph{suekko}) is a native Japanese compound (wago), inherited rather than borrowed from Chinese: \emph{sue}, \textquotedblleft{}end\textquotedblright{} and \emph{ko}, \textquotedblleft{}child\textquotedblright{} (JMdict, KANJIDIC2).
\end{cousinweb}

\begin{grammarlens}[title={What you've built}]
One more word for describing what you lived and felt.
\end{grammarlens}

\subsection*{Your turn}
\begin{itemize}
\raggedright
  \item \emph{Say it:} \emph{suekko}
  \item \emph{Say it:} \emph{suekko}, once more
  \item \emph{Say it:} \emph{suekko wa mada chiisai desu} — The youngest child is still small.
  \item \emph{Recall:} say the Japanese for a parent, then the Japanese for parent and child, then say \emph{suekko} again
\end{itemize}

\begin{tcolorbox}[breakable,colback=teal!4,colframe=teal!35!black,title={Before you move on}]
What does \ja{すえっこ} mean? (\textquotedblleft{}The youngest child\textquotedblright{}.) Say it once more.
\end{tcolorbox}
\section[\texorpdfstring{\ja{めんどう} (mendō)}{めんどう (mendō)}]{\ja{めんどう} (mendō) — a bother; looking after someone}

\label{lesson:JA-C477-mendou}



\begin{quote}
Before the new one: say the Japanese for parent and child, then the Japanese for the youngest child.
\end{quote}

\subsection*{You'll want to know: \ja{めんどう}}
\textbf{\ja{めんどう}} — \emph{mendō} — \textquotedblleft{}a bother; looking after someone\textquotedblright{}.

\begin{quote}
\textbf{\ja{おとうとの} \ja{めんどうを} \ja{みました。}} — \emph{otōto no mendō o mimashita} — \textquotedblleft{}I looked after my younger brother.\textquotedblright{}
\end{quote}

\begin{cousinweb}
\textbf{\ja{めんどう}} (\emph{mendō}) is Sino-Japanese (kango): two Chinese characters read with their on'yomi, readings borrowed from Chinese. They are \emph{men}, \textquotedblleft{}face\textquotedblright{} (Middle Chinese \emph{mjienH}) and \emph{dō}, \textquotedblleft{}collapse\textquotedblright{} (\emph{tawX}) (KANJIDIC2; Middle Chinese in Baxter's transcription).
\end{cousinweb}

\begin{grammarlens}[title={What you've built}]
One more word for describing what you lived and felt.
\end{grammarlens}

\subsection*{Your turn}
\begin{itemize}
\raggedright
  \item \emph{Say it:} \emph{mendō}
  \item \emph{Say it:} \emph{mendō}, once more
  \item \emph{Say it:} \emph{otōto no mendō o mimashita} — I looked after my younger brother.
  \item \emph{Recall:} say the Japanese for parent and child, then the Japanese for the youngest child, then say \emph{mendō} again
\end{itemize}

\begin{tcolorbox}[breakable,colback=teal!4,colframe=teal!35!black,title={Before you move on}]
What does \ja{めんどう} mean? (\textquotedblleft{}A bother; looking after someone\textquotedblright{}.) Say it once more.
\end{tcolorbox}
\section[\texorpdfstring{\ja{しつけ} (shitsuke)}{しつけ (shitsuke)}]{\ja{しつけ} (shitsuke) — upbringing; discipline}

\label{lesson:JA-C477-shitsuke}



\begin{quote}
Before the new one: say the Japanese for the youngest child, then the Japanese for a bother.
\end{quote}

\subsection*{You'll want to know: \ja{しつけ}}
\textbf{\ja{しつけ}} — \emph{shitsuke} — \textquotedblleft{}upbringing; discipline\textquotedblright{}.

\begin{quote}
\textbf{\ja{そふは} \ja{しつけに} \ja{きびしいです。}} — \emph{sofu wa shitsuke ni kibishii desu} — \textquotedblleft{}My grandfather is strict about upbringing.\textquotedblright{}
\end{quote}

\begin{cousinweb}
\textbf{\ja{しつけ}} (\emph{shitsuke}) is a native Japanese word (wago), inherited rather than borrowed from Chinese; its kanji spelling was attached to it later (JMdict).
\end{cousinweb}

\begin{grammarlens}[title={What you've built}]
One more word for describing what you lived and felt.
\end{grammarlens}

\subsection*{Your turn}
\begin{itemize}
\raggedright
  \item \emph{Say it:} \emph{shitsuke}
  \item \emph{Say it:} \emph{shitsuke}, once more
  \item \emph{Say it:} \emph{sofu wa shitsuke ni kibishii desu} — My grandfather is strict about upbringing.
  \item \emph{Recall:} say the Japanese for the youngest child, then the Japanese for a bother, then say \emph{shitsuke} again
\end{itemize}

\begin{tcolorbox}[breakable,colback=teal!4,colframe=teal!35!black,title={Before you move on}]
What does \ja{しつけ} mean? (\textquotedblleft{}Upbringing; discipline\textquotedblright{}.) Say it once more.
\end{tcolorbox}
